%% kapa: exam=1 round=33 year=2022 period=2 form=A subject=appraisal_law subject_ko=감정평가관계법규 seq=5 number=5
%% source: Q-Net 감정평가사 기출문제 (cst00309) artlSeq=5209997 file=raw/1008/5209997/2022년도 제33회 감정평가사 1차시험 2교시(감정평가 관계법규, 회계학).pdf
%% conversion: kordoc 4.24.51 -> tools/gen1.py (auto); review=proofread-sonnet; flags=none
\begin{question}{5}
국토의 계획 및 이용에 관한 법령상 성장관리계획구역을 지정할 수 있는 지역이 아닌 것은? (단, 조례는 고려하지 않음)
\begin{choices}
  \item 개발수요가 많아 무질서한 개발이 진행되고 있거나 진행될 것으로 예상되는 지역
  \item 기반시설이 부족할 것으로 예상되나 기반시설을 설치하기 곤란한 지역을 대상으로 건폐율이나 용적률을 강화하여 적용하기 위한 지역
  \item 「토지이용규제 기본법」 제2조제1호에 따른 지역·지구등의 변경으로 토지이용에 대한 행위제한이 완화되는 지역
  \item 주변의 토지이용이나 교통여건 변화 등으로 향후 시가화가 예상되는 지역
  \item 주변지역과 연계하여 체계적인 관리가 필요한 지역
\end{choices}
\end{question}
\answer{2}
